%# -*- coding: utf-8-unix -*-

\chapter{ELLIPTIC FUNCTIONS}\label{ch:6}

\section{Introduction}

Siegel\index{Siegel, C. L.}\footnote{J.M. 167 (1932), 62--9.} proved in 1932 that if  $\wp(z)$  is a Weierstrass\index{Weierstrass, K.}  $\wp$-function\index{Weierstrass $\wp$-function} such that the invariants  $g_2$, $g_3$ in the equation

\[(\wp'(z))^2 = 4(\wp(z))^3 - g_2\wp(z) - g_3\]
are algebraic numbers\index{algebraic number}, then one at least of any fundamental pair  $\omega$, $\omega'$ of periods of $\wp(z)$ is transcendental; thus both  $\omega$ and $\omega'$ are transcendental if $\wp(z)$ admits complex multiplication\index{complex multiplication}. Siegel\index{Siegel, C. L.}'s work was much improved by Schneider\index{Schneider, Th.}\footnote{M.A. 113 (1937), 1--13.} in 1937; Schneider\index{Schneider, Th.} showed that if $g_2$, $g_3$ are algebraic then any period of $\wp(z)$ is transcendental, and moreover the quotient $\omega/\omega'$ is transcendental except in the case of complex multiplication\index{complex multiplication}. From the latter result it follows at once that the elliptic modular function\index{elliptic modular function} j(z) is transcendental for any algebraic z other than an imaginary quadratic irrational. Schneider\index{Schneider, Th.}'s work led, in fact, to a wide variety of theorems on the transcendence of values of the Weierstrass\index{Weierstrass, K.} functions, and, in 1941, he further obtained far-reaching generalizations concerning Abelian functions\index{Abelian function} and integrals.\footnote{J.M. 183 (1941), 110--28.}

Most of Schneider\index{Schneider, Th.}'s results in this context can be derived as particular cases of a general theorem on meromorphic functions\index{meromorphic function} which he proved in 1949.\footnote{M.A. 121 (1949), 131--40.} The theorem has recently been re-formulated by Lang\index{Lang, S.}.\footnote{See Bibliography (first work).}

\begin{mythm}{}{ch6_1}
Let K be an algebraic number\index{algebraic number} field and let  $f_1(z), \ldots, f_n(z)$ be meromorphic functions\index{meromorphic function} of finite order. Suppose that the ring $K[f_1, \ldots, f_n]$ is mapped into itself by differentiation and has transcendence degree\index{transcendence degree} at least 2 over K. Then there are only finitely many numbers\index{numbers} z at which $f_1, \ldots, f_n$ simultaneously assume values in K.
\end{mythm}

A meromorphic function\index{meromorphic function} f(z) is said to have finite order if there exists $\rho > 0$ and a representation of f as a quotient g/h of entire functions such that, for any $R \ge 2$, and for all z with $|z| \le R$, one has

\[\max(|g(z)|, |h(z)|) < \exp(R^{\rho}). \tag{1}\label{eq:ch6_1}\]

The ring $K[f_1,\dots,f_n]$ consists of all polynomials in $f_1,\dots,f_n$ with coefficients in K, and the transcendence degree\index{transcendence degree} is the maximum number of elements in an algebraically independent subset. \autoref{thm:ch6_1} has been generalized to relate to meromorphic functions\index{meromorphic function} of several variables but the assertion has been obtained only for point sets which can be represented essentially as a cartesian product and this limits considerably the range of application.\footnote{For some work aimed towards overcoming this difficulty see papers by Bombieri\index{Bombieri, E.} (Invent. Math. 10 (1970), 267--87) and Bombieri\index{Bombieri, E.} and Lang\index{Lang, S.} (ibid. 11 (1970), 1--14). It is shown that it suffices if the points in question do not lie on an algebraic hypersurface.} Functions of several variables have been utilized, however, as in Chapters \ref{ch:2} and \ref{ch:3}, in other work on elliptic functions\index{elliptic functions}, and this will be the theme of \S~6.

\section{Corollaries}

We now record some corollaries to \autoref{thm:ch6_1}; others can be found in the works cited in the Bibliography.

\begin{mythm}{}{ch6_2}
If $g_2$, $g_3$ are algebraic, then for any algebraic $\alpha \neq 0$, $\wp(\alpha)$ is transcendental.
\end{mythm}

For the proof one has merely to observe that if $\wp(\alpha)$ were algebraic then, for infinitely many integral values of z, the functions

\[f_1(z) = \wp(\alpha z), \quad f_2(z) = \wp'(\alpha z), \quad f_3(z) = z\]

would simultaneously assume values in the algebraic number\index{algebraic number} field generated by $g_2$, $g_3$, $\alpha$, $\wp(\alpha)$ and $\wp'(\alpha)$ over the rationals, contrary to \autoref{thm:ch6_1}.

\begin{mythm}{}{ch6_3}
For any algebraic $\alpha$ with positive imaginary part, other than a quadratic irrational, $j(\alpha)$ is transcendental.
\end{mythm}

For suppose that $j(\alpha)$ is algebraic. Then there is a $\wp$-function with algebraic invariants $g_2$, $g_3$ and fundamental periods $\omega_1$, $\omega_2$ such that $\alpha = \omega_2/\omega_1$; indeed if $\overline{\wp}(z)$ is the $\wp$-function with periods 1, $\alpha$ and if $\overline{g}_2$, $\overline{g}_3$ are the invariants of $\overline{\wp}$ then the required $\wp$-function has periods $\overline{g}_3^{\frac{1}{3}}$, $\alpha \overline{g}_3^{\frac{1}{3}}$ if $\overline{g}_3 \neq 0$ and $\overline{g}_2^{\frac{1}{2}}$, $\alpha \overline{g}_2^{\frac{1}{2}}$ if $\overline{g}_2 \neq 0$. Now the functions $f_1 = \wp(z)$, $f_2 = \wp(\alpha z)$, $f_3 = \wp'(z)$, $f_4 = \wp'(\alpha z)$ simultaneously assume values in an algebraic number\index{algebraic number} field, say K, when $z = (r + \frac{1}{2})\omega_1$ ($r = 1, 2, \ldots$) and so, by \autoref{thm:ch6_1}, $K[f_1, f_2, f_3, f_4]$ has transcendence degree\index{transcendence degree} at most 1. This implies that $f_1$, $f_2$ are algebraically dependent, whence $l\omega_2$ is a period of $\wp(\alpha z)$ for some positive integer l. Thus $l\alpha \omega_2 = m\omega_1 + n\omega_2$ for some integers m, n and so $\alpha$ is a quadratic irrational. It will be recalled that if 1, $\alpha$ is a basis for an imaginary quadratic field K, then $j(\alpha)$ is in fact a real algebraic integer\index{algebraic integer} with degree given by the class number of K,\footnote{For some work aimed towards overcoming this difficulty see papers by Bombieri\index{Bombieri, E.} (Invent. Math. 10 (1970), 267--87) and Bombieri\index{Bombieri, E.} and Lang\index{Lang, S.} (ibid. 11 (1970), 1--14). It is shown that it suffices if the points in question do not lie on an algebraic hypersurface.} and hence the hypothesis of \autoref{thm:ch6_3} is certainly necessary.

\begin{mythm}{}{ch6_4}
Any vector period of an Abelian function\index{Abelian function} arising from an algebraic curve\footnote{The curve is defined over the algebraic numbers\index{algebraic number}.} by the inversion of Abelian integrals\index{Abelian integral} is transcendental.
\end{mythm}

The result follows from \autoref{thm:ch6_1} with $f_1(z), \dots, f_{n-1}(z)$ given by the Abelian function\index{Abelian function}, say $A(z_1, \dots, z_p)$, and its p partial derivatives with respect to $z_1, \dots, z_p$, evaluated at $z_1 = \omega_1 z, \dots, z_p = \omega_p z$, where $(\omega_1, \dots, \omega_p)$ denotes the given period, together with $f_n(z) = z$. It should perhaps be emphasized that the theorem establishes only the transcendence of one at least of the elements of the period vector, and it remains an open problem to prove the transcendence of each such element. An analogue of \autoref{thm:ch6_4} was obtained by Schneider\index{Schneider, Th.}\footnote{J.M. 183 (1941), 110--28.} in 1941; he showed, by many-variable techniques, that one at least of the first elements of the 2p fundamental periods is transcendental and this implies that the $\beta$-function

\[\beta(a,b) = \int_0^1 x^{a-1} (1-x)^{b-1} dx = \frac{\Gamma(a) \; \Gamma(b)}{\Gamma(a+b)}\]

is transcendental for all rational, non-integral a,b. For if a+b is not an integer then the first elements of the periods of the Abelian function\index{Abelian function} arising from the integration of $x^{a-1}(1-x)^{b-1}$ are given by products of $\beta(a,b)$ with numbers\index{numbers} in the field generated by $e^{2\pi ia}$ and $e^{2\pi ib}$ over the rationals; and the case when a+b is an integer reduces to the transcendence of $\pi\index{transcendence of  $\pi$}$. This result on $\beta(a,b)$ represents all that is known concerning the transcendence of the values of the $\Gamma$-function.

Finally, let $\omega$ be a primitive period of a $\wp$-function with algebraic invariants $g_2$, $g_3$ and let $\eta = 2\zeta(\frac{1}{2}\omega)$ be the associated quasi-period of the Weierstrass\index{Weierstrass, K.} $\zeta$-function\index{Weierstrass $\zeta$-function} satisfying $\zeta'(z) = -\wp(z)$. We have

\begin{mythm}{}{ch6_5}
Any linear combination of $\omega$, $\eta$ with algebraic coefficients, not both 0, is transcendental.
\end{mythm}

For the proof we observe simply that if $\alpha\omega + \beta\eta$ were algebraic, where $\alpha$, $\beta$ are algebraic numbers\index{algebraic number}, not both 0, then the functions

\[f_1 = \wp(z), \quad f_2 = \wp'(z), \quad f_3 = \alpha z + \beta \zeta(z)\]

would simultaneously assume values in an algebraic number\index{algebraic number} field when $z=(r+\frac{1}{2})\,\omega$ ($r=1,2,\dots$), contrary to \autoref{thm:ch6_1}. On recalling that $\omega$ and $\eta$ can be represented as elliptic integrals\index{elliptic integral} of the first and second kinds respectively, one deduces easily from \autoref{thm:ch6_5} that the circumference of any ellipse\index{ellipse} with algebraic axes-lengths is transcendental. Further work in this context will be discussed in \S~6.

\section{Linear equations}\label{sec:ch6_3}

We establish here a result on linear equations\index{linear equations} with algebraic coefficients which generalizes \autoref{lem:ch2_1} of \autoref{ch:2}. K will signify an algebraic number\index{algebraic number} field and $c_1$, $c_2$, $c_3$ will denote positive numbers\index{numbers} that depend on K only. Further, as in \autoref{ch:4}, $\|\theta\|$ will signify the size of $\theta$, that is, the maximum of the absolute values of the conjugates of $\theta$.

\begin{mylemma}{}{ch6_1}
Let M, N be integers with N > M > 0 and let
\[u_{ij} \quad (1 \leqslant i \leqslant M, 1 \leqslant j \leqslant N)\]
be algebraic integers\index{algebraic integer} in K with sizes at most U ($\geqslant 1$). Then there exist algebraic integers\index{algebraic integer} $x_1, \dots, x_N$ in K, not all 0, satisfying
\[\sum_{j=1}^{N} u_{ij} x_j = 0 \quad (1 \leqslant i \leqslant M)\]
and
\[\|x_j\| \le c_1(c_1 N U)^{M/(N-M)} \quad (1 \le j \le N).\]
\end{mylemma}

\begin{proof}
We denote by $\omega_1, \dots, \omega_n$ an integral basis for K and we observe that
\[u_{ij}\omega_k = \sum_{h=1}^n u_{hijk}\omega_h\]
for some rational integers $u_{hijk}$. The equations serve to express the latter as linear combinations of the $u_{ij}$ and their conjugates, with coefficients that depend only on K, and hence we have $|u_{hijk}| < c_2 U$. It follows from \autoref{lem:ch2_1} of \autoref{ch:2} that there exist rational integers $x_{jk}$, not all 0, with absolute values at most $(c_3 NU)^{M/(N-M)}$, satisfying
\[\sum_{j=1}^{N}\sum_{k=1}^{n}u_{hijk}x_{jk}=0 \quad (1\leqslant h\leqslant n,\, 1\leqslant i\leqslant M),\]
and it is now clear that the numbers\index{numbers}
\[x_j = \sum_{k=1}^n x_{jk} \omega_k \quad (1 \leqslant j \leqslant N)\]
have the required properties.
\end{proof}

\section{The auxiliary function}

We assume now that the hypotheses of \autoref{thm:ch6_1} are satisfied and we write $f_i = g_i/h_i$, where $g_i$, $h_i$ are entire functions for which \eqref{eq:ch6_1} holds. We suppose further that there exists a sequence of distinct complex numbers\index{numbers} $y_1, y_2, \dots$ such that $f_i(y_j)$ is an element of K for all i, j. By $c_4, c_5, \dots$ we shall denote positive numbers\index{numbers} which depend only on the quantities so far defined. We signify by m an integer that exceeds a sufficiently large $c_4$, and by k an integer that is sufficiently large compared with m. We write, for brevity, $L = [k^{\frac{1}{2}}]$, and we use $f^{(j)}$ to denote the jth derivative of f.

\begin{mylemma}{}{ch6_2}
There are algebraic integers\index{algebraic integer} $p(\lambda_1, \lambda_2)$ in K, not all 0, with sizes at most $k^{c_2k}$, such that the function
\[\Phi(z) = \sum\limits_{\lambda_1=0}^L \sum\limits_{\lambda_2=0}^L p(\lambda_1,\lambda_2) \left(f_1(z)\right)^{\lambda_1} (f_2(z))^{\lambda_2}\]
satisfies
\[\Phi^{(j)}(y_l)=0 \quad (0\leqslant j\leqslant k,\, 1\leqslant l\leqslant m).\]
\end{mylemma}

\begin{proof}
The number $\Phi^{(j)}(y_l)$ is plainly expressible as a linear form in the $p(\lambda_1, \lambda_2)$ with coefficients given by polynomials in $f_1(y_l), \dots, f_n(y_l)$. The polynomials arise from the derivatives of $f_1, \dots, f_n$ which, by hypothesis, are elements of $K[f_1, \dots, f_n]$; thus the coefficients of $p(\lambda_1, \lambda_2)$ belong to K. The latter become algebraic integers\index{algebraic integer} when multiplied by some positive integer, and we shall suppose that the sizes of these algebraic integers\index{algebraic integer} are at most U. The number of equations to be satisfied is M = m(k+1) and the number of unknowns $p(\lambda_1, \lambda_2)$ is $N = (L+1)^2 > k^{\frac{1}{2}}$. But clearly N > 2M for k sufficiently large and so, by \autoref{lem:ch6_1}, the equations can be solved non-trivially, and indeed with the sizes of the $p(\lambda_1, \lambda_2)$ at most $c_1^2NU$. Hence it remains only to prove that one can take $U \leq k^{c_1k}$.

Now it is readily verified by induction on j that, for any polynomial
\[Q(x_1, \dots, x_n) = \sum_{l_1=0}^{d} \dots \sum_{l_n=0}^{d} q(l_1, \dots, l_n) x_1^{l_1} \dots x_n^{l_n}\]
with coefficients in K, the function $R(z) = Q(f_1, \dots, f_n)$ satisfies
\[R^{(j)}(z) = \sum_{l_1=0}^{d'} \dots \sum_{l_n=0}^{d'} r(l_1, \dots, l_n) f_1^{l_1} \dots f_n^{l_n},\]
where the $r(l_1, \dots, l_n)$ are again elements of K and $d' \leq d + j\delta$, $\delta$ denoting the maximum of the degrees of the first derivatives of $f_1, \dots, f_n$, expressed as polynomials in the latter. Further, it is easily confirmed that if the $q(l_1, \dots, l_n)$ become algebraic integers\index{algebraic integer} with sizes at most s after multiplying Q by some positive integer, then $R^{(j)}$ can be multiplied by a positive integer so that the $r(l_1, \dots, l_n)$ become algebraic integers\index{algebraic integer} with sizes at most $S = (c_2d)^j j! s$. The lemma follows on applying this result with $Q = x_1^{\lambda_1} x_2^{\lambda_2}$ and $j \leq k$, whence $s = 1, d \leq L \leq k$ and $S \leq k^{c_0 k}$, and noting that, if k is sufficiently large, then the estimate $k^{c_0 k}$ obtains for each power product $f_1^{l_1} \dots f_n^{l_n}$ evaluated at $z = y_l$, where $l_i \leq d' \leq c_{10} k$ and $l \leq m$.
\end{proof}

\begin{mylemma}{}{ch6_3}
For any $R \ge 2$ and for all z with $|z| \le R$, the function $\phi = (h_1 \dots h_n)^L \Phi$ satisfies
\[|\phi(z)| < \exp\left\{c_{11}(k\log k + LR^{\rho})\right\}.\]
Further, for any j, l with $j \ge k$, $l \le m$ such that $\Phi^{(i)}(y_l) = 0$ for all i < j, the number $\phi^{(j)}(y_l)$ either vanishes or has absolute value at least $j^{-c_{12}j}$.
\end{mylemma}

\begin{proof}
The first part is an immediate deduction from \eqref{eq:ch6_1} together with the estimates occurring in \autoref{lem:ch6_2}. The second part is obtained by an argument similar to that employed in the proof of \autoref{lem:ch2_3} of \autoref{ch:2}; one observes that $\Phi^{(j)}(y_l)$ is an element of K and that, for $j \ge k$, it becomes an algebraic integer\index{algebraic integer} with size at most $j^{c_{13}j}$ when multiplied by some positive integer likewise bounded. Further, by hypothesis, $\Phi^{(j)}(y_l)$ differs from $\phi^{(j)}(y_l)$ only by a factor $(h_1 \dots h_n)^L$ evaluated at $z = y_l$, and the required result now follows from the fact that the norm of a non-zero algebraic integer\index{algebraic integer} is at least 1.
\end{proof}

\section{Proof of main theorem}

\begin{proof}
It suffices to prove that $\Phi$ vanishes identically; for this implies that $f_1$ and $f_2$ are algebraically dependent and so, since the suffixes can be chosen arbitrarily, $K[f_1, \dots, f_n]$ has transcendence degree\index{transcendence degree} at most 1, contrary to hypothesis. The contradiction shows that m is bounded by some $c_4$ as above, whence the sequence $y_1, y_2, \dots$ must terminate.

The proof will proceed by induction on j; we assume that
\[\Phi^{(i)}(y_l) = 0 \quad (0 \le i < j, \ 1 \le l \le m),\]
and we prove that the same then holds for i = j. In view of \autoref{lem:ch6_2} we can suppose that j > k. Let now C be the circle in the complex plane described in the positive sense with centre the origin and radius $R = j^{1/(4\rho)}$. Further, let
\[F(z) = (z - y_1) \dots (z - y_m),\]
and let l be any integer with $1 \leq l \leq m$. By Cauchy's residue theorem\index{Cauchy's residue theorem}
\[\frac{\phi^{(j)}(y_l)}{(F'(y_l))^j} = \frac{j!}{2\pi i} \int_C \frac{\phi(z)\,dz}{(z-y_l)\,(F(z))^j}.\]

Clearly for z on C we have
\[|F(z)| > (\frac{1}{2}R)^m > j^{m/(8\rho)},\]
and also $|z-y_i| > \frac{1}{2}R$. Further, we have $LR^{\rho} \leq k^{\frac{3}{4}}j^{\frac{1}{4}} \leq j$ and so, by \autoref{lem:ch6_3}, $|\phi(z)| \leq j^{c_{14}j}$. Furthermore, it is obvious that $|F'(y_i)| < j$ for k sufficiently large. Hence we obtain
\[|\phi^{(j)}(y_j)| \leq j^{c_{15}j-jm/(8\rho)}.\]

But if $m > 8\rho(c_{12} + c_{15})$ then, in view of \autoref{lem:ch6_3}, the latter estimate implies that $\phi^{(j)}(y_l) = 0$. Assuming, as plainly one may, that $h_1 \dots h_n$ does not vanish at $z = y_l$, it follows that $\Phi^{(j)}(y_l) = 0$. Thus, by induction, we conclude that $\Phi$ and all its derivatives vanish at $y_1, \dots, y_m$ whence $\Phi$ vanishes identically, as required.
\end{proof}

\section{Periods and quasi-periods}

The work of Siegel\index{Siegel, C. L.}, cited at the beginning, was based on the interpolation techniques discovered a few years previously by Gelfond\index{Gelfond, A. O.},\footnote{See e.g. T{\^o}hoku Math. J. 30 (1929), 280--5.} and the work of Schneider\index{Schneider, Th.} arose out of further developments of these techniques leading, as mentioned in \autoref{ch:2}, to a solution of the seventh problem of Hilbert\index{Hilbert, D.}.\footnote{G{\"{o}}ttingen Nachrichten (1969), No. 16, 145--57.} The recent advances concerning linear forms in the logarithms of algebraic numbers\index{logarithms of algebraic numbers}\index{algebraic number} discussed in earlier chapters have similarly given rise to new results on the transcendental theory of elliptic functions\index{elliptic functions}, as we shall now describe.

First, generalizing \autoref{thm:ch6_5}, it has been shown that if $\omega_1$, $\omega_2$ are primitive periods of some, possibly distinct $\wp$-functions both with algebraic invariants, and if $\eta_1$, $\eta_2$ are the associated quasi-periods of the $\zeta$-functions, we have

\begin{mythm}{}{ch6_6}
Any non-vanishing linear combination of $\omega_1$, $\omega_2$, $\eta_1$, $\eta_2$ with algebraic coefficients is transcendental.
\end{mythm}

This establishes, in particular, the transcendence of the sum of the circumferences of two ellipses\index{ellipse} with algebraic axes-lengths. For the proof of \autoref{thm:ch6_6} we signify by $\wp_1$, $\wp_2$ the given $\wp$-functions, by $\zeta_1$, $\zeta_2$ the associated $\zeta$-functions and we assume, as we may without loss of generality, that the corresponding invariants $\frac{1}{4}g_2$, $\frac{1}{4}g_3$ are algebraic integers\index{algebraic integer}. We assume also that there exists a linear relation
\[\alpha_1\omega_1 + \alpha_2\omega_2 + \beta_1\eta_1 + \beta_2\eta_2 = \alpha_0,\]
where $\alpha_0 \neq 0$, $\alpha_1$, $\alpha_2$, $\beta_1$, $\beta_2$ are algebraic numbers\index{algebraic number}, and we ultimately derive a contradiction. We signify by k an integer which exceeds a sufficiently large number c depending only on the $\alpha$'s, $\beta$'s and the invariants, periods and quasi-periods\index{periods and quasi-periods} of the Weierstrass\index{Weierstrass, K.} functions, and we write, for brevity, $h = [k^{\frac{1}{10}}]$, $L = [k^{\frac{1}{2}}]$. The argument then rests on the construction of an auxiliary function\index{auxiliary function}
\[\Phi(z_1,z_2) = \sum_{\lambda_0=0}^L \sum_{\lambda_1=0}^L \sum_{\lambda_2=0}^L p(\lambda_0,\lambda_1,\lambda_2) \left( f(z_1,z_2) \right)^{\lambda_0} (\wp_1(\omega_1 z_1))^{\lambda_1} (\wp_2(\omega_2 z_2))^{\lambda_2},\]
where the $p(\lambda_0, \lambda_1, \lambda_2)$ are integers, not all 0, with absolute values at most $k^{10k}$, and
\[f(z_1, z_2) = \alpha_1 \omega_1 z_1 + \alpha_2 \omega_2 z_2 + \beta_1 \zeta_1(\omega_1 z_1) + \beta_2 \zeta_2(\omega_2 z_2).\]
The function is constructed to satisfy
\[\Phi_{m_1, m_2}(s+\frac{1}{2}, s+\frac{1}{2}) = 0\]
for all integers s with $1 \le s \le h$ and all non-negative integers $m_1$, $m_2$ with $m_1 + m_2 \le k$, where the suffixes denote partial derivatives as in \autoref{ch:2}.

The essence of the proof is an extrapolation algorithm analogous to that described in connexion with linear forms in logarithms\index{linear forms in logarithms}, but the order of $\Phi$ here is greater than in the earlier work and, to compensate, rational extrapolation points with large denominators are utilized; an important role in the discussion is therefore played by the division value properties of the elliptic functions\index{elliptic functions}. The counterpart of Lemma 4 of \autoref{ch:2} asserts that, for any integer J between 0 and 50 inclusive, we have
\[\Phi_{m_1, m_2}(s + r/q, s + r/q) = 0\]
for all integers q, r, s with q even, (r, q) = 1,
\[1\leqslant q\leqslant 2h^{\frac{1}{4}J},\quad 1\leqslant s\leqslant h^{\frac{1}{4}J+1},\quad 1\leqslant r< q,\]
and all non-negative integers $m_1$, $m_2$ with $m_1 + m_2 \leq k/2^J$. The demonstration proceeds by induction and involves an application of the maximum-modulus principle\index{maximum-modulus principle} as in the original lemma. It also utilizes the observation that, apart from a factor $\omega_1^{m_1}\omega_2^{m_2}$, the number on the left of the required equation is algebraic with degree at most $c'q^4$, where c' is defined like c above; and precise estimates for the number and its conjugates are furnished by division value theory. One concludes from the lemma that
\[\Phi_{m_1, m_2}(s+\frac{1}{4}, s+\frac{1}{4}) = 0 \quad (1 \le s \le L+1, \ 0 \le m_1, m_2 \le L),\]
which is clearly a system of $(L+1)^3$ linear equations\index{linear equations} in the same number of variables $p(\lambda_0, \lambda_1, \lambda_2)$; on noting that, for any regular function f, the determinant or order n with the ith derivative of $(f(z))^j$ in the ith row and jth column has value
\[2! \dots n! (f'(z))^{\frac{1}{2}n(n+1)},\]
one easily verifies that the system of equations is untenable, and this proves \autoref{thm:ch6_6}.

The special case of the theorem when $\wp_1$, $\wp_2$ are the same $\wp$-function, say $\wp$, is of particular interest. For then $\omega_1$, $\omega_2$ can be taken as a pair of fundamental periods of $\wp$ and we have the Legendre relation\index{Legendre relation}
\[\eta_1\omega_2-\eta_2\omega_1=2\pi i.\]
In this case Coates\index{Coates, J.}\footnote{Amer. J. Math. 93 (1971), 385--97; Inventiones Math. 11 (1970), 167--82.} and more recently Masser\index{Masser, D. W.}\footnote{Ph.D. Thesis, Cambridge, 1974.} have much extended the arguments and have proved:

\begin{mythm}{}{ch6_7}
The space spanned by 1, $\omega_1$, $\omega_2$, $\eta_1$, $\eta_2$ and $2\pi i$ over the algebraic numbers\index{algebraic number} has dimension either 4 or 6 according as $\wp$ does or does not admit complex multiplication\index{complex multiplication}.
\end{mythm}

The theorem clearly exhibits a non-trivial example of five numbers\index{numbers} that are algebraically dependent but linearly independent over the algebraic numbers\index{algebraic number}. Moreover it implies that, when $\wp$ admits complex multiplication\index{complex multiplication}, the numbers in question satisfy an algebraic linear relation other than that between the periods; this was discovered by Masser\index{Masser, D. W.}. It takes the form
\[a\eta_{2}-c\tau\eta_{1}=\gamma\omega_{1}\]
where $\gamma$ is algebraic and a, c are the integers occurring in the equation
\[a + b\tau + c\tau^2 = 0\]
satisfied by $\tau = \omega_1/\omega_2$. A necessary and sufficient condition for $\gamma$ to be 0 is that either $g_2$ or $g_3$ be 0, and thus one deduces that $\eta_1/\eta_2$ is transcendental if and only if neither invariant vanishes. The theorem also shows, for instance, that $\pi + \omega$ and  $\pi + \eta$ are transcendental for any period $\omega$ of $\wp(z)$ and quasi-period $\eta$ of $\wp(z)$. The transcendence of $\pi\index{transcendence of  $\pi$}$, incidentally, follows from \autoref{thm:ch6_1} by way of the functions $\wp(\omega z/\pi)$ and $e^{2\pi iz}$.

The demonstration of \autoref{thm:ch6_6} extends easily to establish, under the conditions appertaining to \autoref{thm:ch6_7}, the transcendence of any non-vanishing linear combination of $\omega_1$, $\omega_2$, $\eta_1$, $\eta_2$ and $2\pi i$; the auxiliary function\index{auxiliary function} now takes the form
\[ \Phi(z_1,z_2,z_3) = \sum_{\lambda_0=0}^L \dots \sum_{\lambda_3=0}^L p(\lambda_0,\dots,\lambda_3) (f(z_1,z_2,z_3))^{\lambda_0} (\wp_1(\omega_1z_1))^{\lambda_1} (\wp_2(\omega_2z_2))^{\lambda_2} e^{2\pi i \lambda_3 z_3}, \]
where $L = [k^{\frac{1}{2}}]$ and $f(z_1, z_2, z_3)$ is the sum of $f(z_1, z_2)$, as defined above, and an algebraic multiple of $\pi z_3$. Here, however, it is necessary to appeal to another remarkable property of the division values, namely that, for any positive integer n, the field obtained by adjoining $\wp(\omega_1/n)$, $\wp(\omega_2/n)$, $\wp'(\omega_1/n)$ and $\wp'(\omega_2/n)$ to $K = \mathbb{Q}(g_2, g_3, e^{2\pi i/n})$ has degree at most $2n^3$ over K; this ensures that the estimate $c'q^4$ referred to above remains unaltered in the present context. To complete the proof of \autoref{thm:ch6_7} one has to establish the linear independence over the algebraic numbers\index{algebraic number} of $\omega_1$, $\eta_1$ and $2\pi i$ in the case when $\wp$ admits complex multiplication\index{complex multiplication}, and of these, together with $\omega_2$, $\eta_2$, in the case when $\wp$ does not. The work runs on similar lines, using slightly modified auxiliary functions\index{auxiliary function}, but the determinant arguments at the end are no longer applicable; ad hoc techniques have been introduced to overcome this difficulty involving, in particular, new considerations on the density of zeros of meromorphic functions\index{meromorphic function}. The linear independence of $\omega_1, \omega_2$ and $2\pi i$ was in fact proved first by Coates\index{Coates, J.} utilizing a deep result of Serre\index{Serre, J. P.}, but Masser\index{Masser, D. W.} later verified this more elementarily.

In another direction, the work has been refined to yield estimates from below for linear forms in periods and quasi-periods\index{periods and quasi-periods}. They show, for instance, that for any $\wp$-function with algebraic invariants, for any $\varepsilon > 0$, and for any positive integer n,
\[|\wp(n)| < C n^{(\log \log n)^{7+\epsilon}},\]
where C depends only on $g_2$, $g_3$ and $\epsilon$.\footnote{Amer. J. Math. 92 (1970), 619--22 (A. Baker\index{Baker, R. C.}\index{Baker, A.}); P.C.P.S. 73 (1973), 339--50 (D. W. Masser\index{Masser, D. W.}).} In fact a similar result has been established for $\wp(\pi+n)$ and for $\wp(\alpha)$, where $\alpha$ is any non-zero algebraic number\index{algebraic number}. The estimate compares well with the lower bound $|\wp(n)| > Cn$ valid for some C > 0 and infinitely many n.

Finally, as a further example of the type of theorem that has been obtained by the above methods, we mention a recent result of Masser\index{Masser, D. W.}\footnote{Ph.D. Thesis, Cambridge, 1974.} concerning algebraic points on elliptic curves\index{elliptic curve}; he has proved, namely, that if $\wp(z)$ has algebraic invariants and admits complex multiplication\index{complex multiplication}, then any numbers\index{numbers} $u_1, \ldots, u_n$ for which $\wp(u_i)$ is algebraic are either linearly dependent over $\mathbb{Q}(\tau)$ or linearly independent over the field of all algebraic numbers\index{algebraic number}. It would be of much interest to establish a theorem of the latter kind more generally for all $\wp$-functions with algebraic invariants, and it would likewise be of interest to extend \autoref{thm:ch6_6} to apply to any number of $\wp$-functions; both problems, however, seem out of reach at present.
